\documentclass[11pt]{article}
\usepackage{amsmath,amssymb,graphicx,bm}
\usepackage[a4paper,margin=1in]{geometry}

\title{\textbf{The Endogenized Intentional Stance:\\
A Predictive–Compression Formalization of Agent Discovery}}
\author{Author Name\\AE Studio\\\texttt{author@aestudio.org}}
\date{November 2025}

\begin{document}
\maketitle

\begin{abstract}
We formalize an \emph{Endogenized Intentional Stance (EIS)}: a behavioral–informational procedure for discovering agents and goals directly from raw observation streams without pre-specified entities or sensors.  EIS extends the pragmatic–predictive view of Dennett (1987) and recent stochastic formalizations (McGregor et al., 2025) by replacing the assumption of a given world model with a learned latent dynamics that simultaneously grounds environment, action, and purpose.  Agency emerges where a soft-optimal (goal-rational) prior over latent interventions reduces description length relative to a purely mechanistic baseline.  Intentionality thus becomes an empirical, compression-based property of dynamics rather than an anthropomorphic stance.  We derive sufficient conditions for intentional gain, outline a hierarchical aggregation scheme, and connect the resulting constructs to known information-theoretic quantities: predictive mutual information, control entropy, and value gradients under bounded rationality.
\end{abstract}

\section{Introduction}

Attributing beliefs and desires to a system can dramatically simplify behavioral prediction.  Dennett’s \emph{intentional stance} treats such attributions as modeling conveniences, vindicated when they increase predictive success at low cognitive cost.  McGregor et al. (2025) showed that any stochastic process may, in principle, admit a normative–epistemic interpretation.  Yet both perspectives presuppose a \emph{world model} defining what could be believed or desired.  We here remove that presupposition: given only an observation sequence $\{o_t\}_{t=1}^T$, we ask where goal-directed explanations yield better compression than purely mechanistic ones.  The resulting procedure endogenizes world, agent, and goal in a single generative model.

\section{Predictive world modeling}

Let $\theta$ parameterize a predictive latent dynamics
\begin{equation}
p_\theta(o_{1:T},z_{1:T})=
\prod_{t} p_\theta(o_t\mid z_t)\,p_\theta(z_t\mid z_{t-1}),
\label{eq:latent}
\end{equation}
with hidden variables $z_t\in\mathbb R^d$ learned by maximizing
\[
\mathcal{L}_0(\theta,\phi)
=\sum_t \mathbb E_{q_\phi(z_t\mid o_{\le t})}
[\log p_\theta(o_t\mid z_t)]-\beta\,D_{\mathrm{KL}}(q_\phi\Vert p_\theta).
\]
No partition into sensors, actuators, or entities is assumed.

\section{Intentional prior and compression gain}

We augment~(\ref{eq:latent}) with discrete latent impulses
$a_t\in\mathcal A$ acting as candidate interventions in $z$:
\begin{equation}
p(o,z,a)=
\prod_t p(o_t\mid z_t)\,p(z_{t+1}\mid z_t,a_t)\,p(a_t).
\end{equation}
The \emph{intentional stance} introduces a soft-optimal prior
\begin{equation}
p_\beta(a_t\mid z_t)
\propto \exp[\beta\,r_\psi(z_t,a_t)],
\label{eq:soft}
\end{equation}
where $r_\psi$ is a learned reward density and $\beta>0$ an inverse-temperature.
This converts the purely generative dynamics into a bounded-rational control process.

Let $\mathcal{L}_\mathrm{I}$ denote the log evidence of the intentional model and
$\mathcal{L}_\mathrm{M}$ that of a baseline with uniform $p(a_t)$.
The \emph{intentional gain}
\begin{equation}
\Delta\mathcal L
=\mathcal{L}_\mathrm{I}-\mathcal{L}_\mathrm{M}
-\lambda\,\mathrm{DL}(r_\psi)
\label{eq:gain}
\end{equation}
measures the code-length reduction due to goal-directed explanation.
A subsystem $C$ is declared \emph{intentional} when $\Delta\mathcal L_C>0$,
signifying that rationality assumptions compress its trajectory.

\section{Emergent entities}

We seek subsets of latent coordinates $C\subseteq z$ for which $\Delta\mathcal L_C$ is maximized subject to sparse coupling with the remainder.
Operationally, we minimize
\begin{equation}
\mathcal{J}(C)=
I(z'_{C};z'_{\neg C}\mid a_C)
-\alpha\,\Delta\mathcal L_C,
\end{equation}
where $I(\cdot;\cdot\mid\cdot)$ is conditional mutual information.
Local minima define dynamically quasi-autonomous modules whose behavior is efficiently summarized by~(\ref{eq:soft}).
These modules constitute \emph{endogenously discovered agents}.

\section{Hierarchical aggregation}

Given intentional modules $\{C_i\}$ with latent values $z^{(i)}_t$,
construct an abstract state $\bar z_t=f(\{z^{(i)}_t\})$ and abstract actions $\bar a_t=g(\{a^{(i)}_t\})$.
Reapplying~(\ref{eq:soft}) in the compressed space yields higher-level intentional priors
$p_\beta(\bar a_t\mid\bar z_t)\!\propto\!\exp[\beta\,\bar r(\bar z_t,\bar a_t)]$.
Hierarchies thus emerge recursively where joint intentional gain exceeds the sum of constituent gains.
This realizes Dennett’s \emph{multi-level stance}: agents within agents, unified by shared compression criteria.

\section{Cooperative structure}

For modules $i,j$ define an empirical \emph{cooperativity index}
\begin{equation}
\kappa_{ij}
=\frac{\Delta V^{(i)}( \text{do}(a^{(j)}) )}
      {\mathrm{cost}^{(i)}},
\label{eq:kappa}
\end{equation}
where $\Delta V^{(i)}$ is the change in expected cumulative reward
$V^{(i)}(z)=\mathbb E[\sum_t r_\psi^{(i)}(z_t,a_t)]$
under counterfactual manipulation of $a^{(j)}$.
Edges with $\kappa_{ij}>1$ define cooperative clusters;
$\kappa_{ij}<0$ indicates competition.
These relations permit percolation analyses analogous to those used in evolutionary models of cooperation
(Hamilton, 1964).

\section{Relation to information principles}

The intentional prior~(\ref{eq:soft}) connects to known information-theoretic quantities:
let $I_{\mathrm{pred}}=I(z_t;z_{t+1})$ and
$I_{\mathrm{ctrl}}=I(a_t;z_{t+1})$.
Expected log evidence decomposes as
\begin{equation}
\mathbb E[\log p(o)]\approx
I_{\mathrm{pred}}+I_{\mathrm{ctrl}}
-\beta H(z_t)-\gamma S,
\end{equation}
with $S=-\mathbb E[\log p(o_{t+1}\mid z_t)]$ the residual surprise
and coefficients $(\beta,\gamma)$ reflecting memory and uncertainty costs
(Friston, 2010; Tishby et al., 2000).
Hence intentional gain~(\ref{eq:gain}) corresponds to an information-bottleneck advantage achieved by positing latent goals.

\section{Determinacy and uniqueness}

Following McGregor et al. (2025),
multiple normative–epistemic interpretations can coexist for stochastic processes.
EIS refines this by assigning quantitative preference
\[
P(\text{interpretation})\propto
\exp[\Delta\mathcal L],
\]
so that interpretations form an energy landscape whose minima represent the simplest rationalizations.
Deterministic subsystems yield unique minima;
chaotic or noisy ones support degenerate stances.

\section{Discussion}

EIS reinterprets Dennett’s intentional stance as a model-selection principle.
Whereas classic formulations required an observer’s external world model,
EIS internalizes it by training a predictive generative model directly from data.
Intentionality becomes an empirical regularity—regions of the dynamics for which
rational-goal assumptions save bits.
This extends bounded-rational control to settings lacking predefined agents or utilities,
aligning with information-theoretic frameworks of adaptation
(Friston et al., 2010; Tishby et al., 2000)
and with normative–epistemic mappings in stochastic processes
(McGregor et al., 2025).

\section{Conclusion}

The Endogenized Intentional Stance provides a unified,
data-driven criterion for the emergence of agency and purpose.
An observer need not know what is sensed or controlled;
the system’s own predictive structure reveals where such attributions
improve compression.
Hierarchies and cooperation follow naturally from
additivity of intentional gain.
This formulation renders the intentional stance measurable,
bridging philosophical pragmatism and information physics.

\begin{thebibliography}{9}

\bibitem{Dennett1987}
D.~C.~Dennett.
\newblock \emph{The Intentional Stance}.
\newblock MIT Press, 1987.

\bibitem{Friston2010}
K.~Friston.
\newblock The free-energy principle: a unified brain theory?
\newblock \emph{Nature Reviews Neuroscience}, 11(2):127–138, 2010.

\bibitem{Hamilton1964}
W.~D.~Hamilton.
\newblock The genetical evolution of social behaviour.
\newblock \emph{Journal of Theoretical Biology}, 7(1):1–16, 1964.

\bibitem{Tishby2000}
N.~Tishby, F.~C.~Pereira, and W.~Bialek.
\newblock The information bottleneck method.
\newblock \emph{arXiv:physics/0004057}, 2000.

\bibitem{McGregor2025}
S.~McGregor, et al.
\newblock Formalising the intentional stance I: Attributing goals and beliefs to stochastic processes.
\newblock \emph{arXiv:2405.16490v2}, 2025.

\end{thebibliography}

\end{document}

